% !PRINTED 176
\weberpage{176}
%<<CONT
sind. Es findet daher (mit Rücksicht auf \S\,45) eine der folgenden
Beziehungen statt:
\[
\lambda^2 = \varkappa^2(\omega), \quad \varkappa^2\!\left(-\frac{1}{\omega}\right), \quad \varkappa^2(\omega+1), \quad \varkappa^2\!\left(\frac{\omega-1}{\omega}\right),
\]
\[
\varkappa^2\!\left(-\frac{1}{\omega+1}\right), \quad \varkappa^2\!\left(\frac{\omega}{1-\omega}\right),
\]
wonach aus dem ersten Satze die Richtigkeit des zweiten folgt.

Die Variable $\omega$ ist hier immer auf einen positiven imaginären
Teil beschränkt, und wenn wir zwei durch eine Gleichung (11)
zusammenhängende Zahlen $\omega$, $\omega'$ \gesperrt{äquivalent} nennen, so haben
alle mit $\omega$ äquivalenten Zahlen einen positiven imaginären Teil.
Wir können unseren Satz 2. auch so aussprechen:

3. \gesperrt{Die Funktion $j(\omega)$ hat für alle äquivalenten Werte $\omega$
und nur für diese einen und denselben Wert}$^{1)}$.

Zu jedem Wert von $\omega$ gehört ein bestimmter Wert von $j(\omega)$,
und zu jedem Wert von $j(\omega)$ ein Wert von $\omega$ und die Gesamtheit
 der äquivalenten Werte $\omega'$. Wenn also eine einwertige
Funktion $\Phi(\omega)$ von $\omega$ der Bedingung genügt:
\[
\Phi\!\left(\frac{\gamma+\delta\omega}{\alpha+\beta\omega}\right) = \Phi(\omega),
\]
so ist $\Phi(\omega)$ eine einwertige Funktion von $j(\omega)$, und wenn wir
also annehmen, daß $\Phi(\omega)$ nur für eine endliche Anzahl von
Werten $j(\omega)$ unendlich und nur in endlicher Ordnung unendlich
werde, so folgt, daß $\Phi(\omega)$ eine \gesperrt{rationale Funktion von $j(\omega)$} ist.

Die Voraussetzung trifft z.\,B. immer dann zu, wenn $\Phi(\omega)$
mit $j(\omega)$ in einem \gesperrt{algebraischen Zusammenhange} steht.

\begin{center}
\textbf{\S\ 54. \gesperrt{Die Modulfunktionen.}}
\end{center}

Es sei $\psi(\omega)$ eine eindeutige Funktion von $\omega$. Wendet man
auf $\omega$ eine lineare Transformation $S = \begin{pmatrix} \alpha, & \beta \\ \gamma, & \delta \end{pmatrix}$ an, so geht $\psi(\omega)$
in eine andere Funktion
\weq{1}{\psi\!\left(\frac{\gamma+\delta\omega}{\alpha+\beta\omega}\right)}
über, die wir mit $\psi\,|\,S$ bezeichnen wollen.

\begingroup
\renewcommand{\thefootnote}{1)}
\footnotetext[1]{Dieser Satz ist zuerst von \gesperrt{Dedekind} bewiesen \gesperrt{(Crelle, Bd.\,83)}.
\gesperrt{Dedekind} nennt danach die Funktion $\mathrm{val}(\omega) = 3^{-3}2^{-6}j(\omega)$ die \gesperrt{Valenz}
von $\omega$.}
\endgroup
